% !TEX encoding = UTF-8 Unicode
% !TEX TS-program = pdflatex
% !TEX root = ../Tesi.tex



%************************************************
\myChapter{Problemi al contorno}
%************************************************
\label{cap:bvp}

In questo capitolo viene presentata la formulazione di un problema al contorno consistente (\textit{Boundary Value Problem, BVP})  per le equazioni differenziali derivate dal problema elettromagnetico presentato nel Capitolo \ref{cap:maxwell}. 


I potenziali introdotti in Sezione \ref{subsec:1potenziale} ben si prestano alla scrittura in forma debole delle equazioni di Maxwell. In questo modo, con riferimento al problema originario, le due equazioni \eqref{eq:2max1} e \eqref{eq:2max4} saranno soddisfatte in forma forte.

Verrà inoltre da qui in avanti abbandonato il regime dinamico completo per focalizzare l'attenzione sui regimi ridotti adeguati alle analisi effettuate in questo lavoro.

Infine, verranno utilizzati unicamente legami costitutivi isotropi, così da ridurre i tensori $\TT{\varepsilon}$ e $\TT{\mu}$ alle quantità scalari $\varepsilon$ e $\mu$.

Come si è potuto osservare nel capitolo precedente, nelle due formulazioni ridotte \textit{MQS} e stazionaria le equazioni sono disaccoppiate nei due campi $\left\{\T{a},\phi\right\}$:
\begin{equation}
\left\{
\label{eq:3potqss}
\begin{aligned}
\lap\left({\phi}\right) &= -\varepsilon^{-1} \rho\\
-\lap\left(\T{a}\right) &= \mu\T{j}+\rot\left(\T{b_r}\right).
\end{aligned}
\right.
\end{equation}
I due regimi si differenziano nel recupero dei campi $\left\{\T{e},\T{b}\right\}$:

\begin{equation}
\label{eq:2recuperoqss}
\begin{aligned}
&\text{MQS}\\
&\left\{
\begin{aligned}
\T{b}&=\rot \left(\T{a}\right)\\
\T{e} &= -\grad\left(\phi\right) -\partial_t \T{a}
\end{aligned}
\right.
\end{aligned}
\qquad
\begin{aligned}
&\text{STAZIONARIO}\\
&\left\{
\begin{aligned}
\T{b}&=\rot \left(\T{a}\right)\\
\T{e} &= -\grad\left(\phi\right)
\end{aligned}
\right.
\end{aligned}
\end{equation}

Durante la presentazione dei potenziali si è evidenziato l'utilizzo per i regimi ridotti della condizione di gabbia di Gauss \eqref{eq:2gaussgauge}. L'imposizione di tale condizione comporta l'interscambiabilità dei seguenti termini:
\begin{equation}
\label{eq:3rotrotlap}
\rot\left[\rot\left(\T{a}\right)\right] = -\lap\left(\T{a}\right)
\end{equation}
all'interno della seconda delle equazioni \eqref{eq:3potqss}. 

Verranno ora introdotti gli ambienti naturali per lo sviluppo della forma debole, per poi analizzare in dettaglio le formulazioni per il problema magnetico ed elettrico.

\section{Spazi funzionali}

Si definisca lo spazio
\begin{equation}
\label{eq:3L2}
L_2\left(\TTTT{D}\right) := \left\{ u | \int_{\TTTT{D}} u^2 dv < \infty\right\},
\end{equation}
che fornito del prodotto scalare
\begin{equation}
\label{eq:3scaL2}
\left(u,v\right)= \int_{\TTTT{D}} u\,v\, dv,
\end{equation}
costituisce uno spazio di Hilbert.

Si richiamano ora le definizioni degli \textit{operatori differenziali generalizzati}:
\begin{enumerate}


\item $D_u = u'$ è detta \textit{derivata generalizzata} di $u \in L_2\left(\TTTT{D}\right)$ se 
\begin{equation}
\label{eq:3gen_der}
\int_{\TTTT{D}}D_u w\, dv = - \int_{\TTTT{D}}u \,w' \, dv \qquad \forall w \in C_0^\infty\left(\bar{\TTTT{D}}\right),
\end{equation}
dove $\bar{\TTTT{D}}$ è un sottoinsieme compatto di $\TTTT{D}$ e costituisce il supporto delle \textit{test function} $w$. 
\item $\T{g}=\grad\left(u\right)$ è detto \textit{gradiente generalizzato} di $u \in  L_2\left(\TTTT{D}\right)$ se
\begin{equation}
\label{eq:3gen_grad}
\int_{\TTTT{D}}\T{g}\cdot \T{w}\, dv = - \int_{\TTTT{D}}u \,\dive\left(\T{w}\right) \, dv \qquad \forall \T{w} \in \left(C_0^\infty\left(\bar{\TTTT{D}}\right)\right)^3.
\end{equation}
\item $\T{r}=\rot\left(\T{u}\right)$ è detto \textit{rotore generalizzato} di $\T{u} \in  \left(L_2\left(\TTTT{D}\right)\right)^3$ se
\begin{equation}
\label{eq:3gen_rot}
\int_{\TTTT{D}}\T{r}\cdot \T{w}\, dv =  \int_{\TTTT{D}}\T{u}\cdot \rot\left(\T{w}\right) \, dv \qquad \forall \T{w} \in \left(C_0^\infty\left(\bar{\TTTT{D}}\right)\right)^3.
\end{equation}
\item $d = \dive\left(\T{u}\right)$ è detta \textit{divergenza generalizzata} di $\T{u} \in  \left(L_2\left(\TTTT{D}\right)\right)^3$ se
\begin{equation}
\label{eq:3gen_div}
\int_{\TTTT{D}}d\, w\, dv =  \int_{\TTTT{D}}\T{u}\cdot \grad\left(w\right) \, dv \qquad \forall w \in C_0^\infty\left(\bar{\TTTT{D}}\right).
\end{equation}
\end{enumerate}

Si possono ora definire i seguenti spazi funzionali:
\begin{eqnarray}
\label{eq:3h1}
H^1\left(\TTTT{D}\right) &:=&\left\{ w \in L_2\left(\TTTT{D}\right) | \grad\left( w\right) \in \left(L_2\left(\TTTT{D}\right)\right)^3\right\},\\
\label{eq:3hcurl}
H\left(\rot,\TTTT{D}\right) &:=&\left\{ \T{w} \in \left(L_2\left(\TTTT{D}\right)\right)^3 | \rot\left( \T{w}\right) \in \left(L_2\left(\TTTT{D}\right)\right)^3\right\},\\
\label{eq:3hdiv}
H\left(\dive,\TTTT{D}\right) &:=&\left\{ \T{w} \in \left(L_2\left(\TTTT{D}\right)\right)^3 | \dive\left( \T{w}\right) \in L_2\left(\TTTT{D}\right)\right\}.
\end{eqnarray}

A questi vengono associati i seguenti prodotti scalari:
\begin{eqnarray}
\label{eq:3scah1}
\left(u,v\right)_1 &=& \int_{\TTTT{D}} \grad\left(u\right)\cdot\grad\left(v\right)dv+\left(u,v\right),\\
\label{eq:3scahcurl}
\left(\T{u},\T{v}\right)_{\rot} &=& \int_{\TTTT{D}}\rot\left(\T{u}\right)\cdot\rot\left(\T{v}\right)dv+\left(\T{u},\T{v}\right),\\
\label{eq:3scahdiv}
\left(\T{u},\T{v}\right)_{\dive} &=& \int_{\TTTT{D}}\dive\left(\T{u}\right)\dive\left(\T{v}\right)dv+\left(\T{u},\T{v}\right).
\end{eqnarray}

Gli spazi sopra definiti, se forniti delle norme associate ai rispettivi prodotti scalari,  sono spazi di Hilbert.

\section{Forma debole per il problema magnetico}

Si consideri l'equazione:
\begin{equation}
\rot\left[\rot\left(\T{a}\right)\right] = \mu\T{j}+\rot\left(\T{b_r}\right),
\label{eq:3strong1}
\end{equation}
alla quale vengono associate le condizioni di interfaccia derivanti dalle equazioni \eqref{eq:4interface_b3} e \eqref{eq:4interface_h}:
\begin{equation}
\left\{
\begin{aligned}
&\Bigl\langle \rot\left(\T{a}\right)\times \T{n}\Bigr\rangle = -\mu\T{j_s} \\
&\Bigl\langle \rot\left(\T{a}\right)\cdot \T{n}\Bigr\rangle = 0.
\end{aligned}
\right.
\label{eq:3interface_a}
\end{equation}

Si noti che la seconda delle due condizioni al contorno è equivalente a
\begin{equation}
\label{eq:3interface_a2}
\Bigl\langle \T{a}\times \T{n}\Bigr\rangle = 0.
\end{equation}
\\
Sia ora $\T{a}\in H\left(\rot,\TTTT{D}\right)$; per ogni dominio $\TTTT{D}$ si integri l'equazione \eqref{eq:3strong1} e la prima delle condizioni al contorno \eqref{eq:3interface_a} e le si moltiplichino per una \textit{test function} vettoriale $\T{w} \in H\left(\rot,\TTTT{D}\right)$:
\begin{align}
\label{eq:3weak1}
&\int_{\TTTT{D}}\T{w} \cdot \rot\left[\rot\left(\T{a}\right)\right] dv+\int_{\partial\TTTT{D}}\T{w}\cdot \left[\rot\left(\T{a}\right)\times\T{n}\right]ds= \nonumber\\
&=\int_{\TTTT{D}}\T{w}\cdot\left[\mu\T{j}+rot\left(\T{b_r}\right)\right]dv+\int_{\partial\TTTT{D}}\T{w}\cdot \mu\T{j_s} ds.
\end{align}

Si sfrutti ora la relazione vettoriale
\begin{equation}
\dive\left[\T{u}\times \rot\left(\T{v}\right)\right] = \rot\left(\T{u}\right)\cdot \rot\left(\T{v}\right)-\T{u}\cdot \rot\left[rot\left(\T{v}\right)\right]
\label{eq:3vectrel}
\end{equation}
per riscrivere l'equazione \eqref{eq:3weak1}:
\begin{align}
\label{eq:3weak2}
&\int_{\TTTT{D}}\rot\left(\T{w}\right)\cdot \rot\left(\T{a}\right) dv-\int_{\TTTT{D}} \dive\left[\T{w}\times \rot\left(\T{a}\right)\right] dv + \nonumber\\
&+\int_{\partial\TTTT{D}}\T{w}\cdot \left[\rot\left(\T{a}\right)\times\T{n}\right]ds= \nonumber\\
&=\int_{\TTTT{D}}\T{w}\cdot\left[\mu\T{j}+\rot\left(\T{b_r}\right)\right]dv+\int_{\partial\TTTT{D}}\T{w}\cdot \mu\T{j_s} ds.
\end{align}

Si applichi ora il teorema di Gauss al termine di divergenza:
\begin{align}
\label{eq:3weak3}
&\int_{\TTTT{D}}\rot\left(\T{w}\right)\cdot \rot\left(\T{a}\right) dv-\int_{\partial\TTTT{D}} \left[\T{w}\times \rot\left(\T{a}\right)\right]\cdot \T{n} ds + \nonumber\\
&+\int_{\partial\TTTT{D}}\T{w}\cdot \left[\rot\left(\T{a}\right)\times\T{n}\right]ds= \nonumber\\
&=\int_{\TTTT{D}}\T{w}\cdot\left[\mu\T{j}+\rot\left(\T{b_r}\right)\right]dv+\int_{\partial\TTTT{D}}\T{w}\cdot \mu\T{j_s} ds.
\end{align}

Si può dimostrare che i due termini in $\partial\TTTT{D}$ nel primo membro dell'Eq. \eqref{eq:3weak2} sono uguali citando l'identità vettoriale:
\begin{equation}
\label{eq:3vectidentity}
\T{p} \cdot (\T{q} \times \T{r}) = \T{q} \cdot (\T{r} \times \T{p}) = \T{r} \cdot (\T{p} \times \T{q}).
\end{equation}

Si ottiene pertanto:
$\forall \,\T{w} \in H(\rot,D)$:
\begin{align}
\label{eq:3weak4}
&\int_{\TTTT{D}}\rot\left(\T{w}\right)\cdot \rot\left(\T{a}\right) dv
= \nonumber\\
&=\int_{\TTTT{D}}\T{w}\cdot\left[\mu\T{j}+\rot\left(\T{b_r}\right]\right)dv+\int_{\partial\TTTT{D}}\T{w}\cdot \mu\T{j_s} ds.
\end{align}

A questa forma debole deve essere incorporata la condizione di gabbia sul potenziale \eqref{eq:2gaussgauge} unitamente alla seconda delle condizioni al contorno \eqref{eq:3interface_a}; si introducano una \textit{test function}  $\psi$ e la corrispondente \textit{trial function} $\chi$ appartenenti allo spazio funzionale $H^1\left(\TTTT{D}\right)$. Si ha, sfruttando la definizione di divergenza generalizzata \eqref{eq:3gen_div}, l'ortogonalità del potenziale vettoriale rispetto ai gradienti:
\begin{equation}
\label{eq:3weakgauge}
\int_{\TTTT{D}} \T{a} \cdot \grad\left(\psi\right) dv = 0 \qquad \forall \psi \in H^1\left(\TTTT{D}\right).
\end{equation}

 Si ottiene così la cosiddetta \textit{formulazione magnetostatica mista}, per la cui soluzione esistenza e unicità possono essere dimostrate mediante il Teorema di \textit{Brezzi} e il Teorema di \textit{Lax-Milgram}, omessi in questa trattazione \citep{zaglmayr_high_2006}.
 
\begin{itemize}
\item $\forall \T{w} \in H(\rot,\TTTT{D})$:
\begin{align}
\label{eq:3weak3_a}
\int_{\TTTT{D}_i}\rot\left(\T{w}\right)\cdot \rot\left(\T{a}\right) dv +\int_{\TTTT{D}_i} \T{w} \cdot \grad \left(\chi\right) dv
= \nonumber\\
=\int_{\TTTT{D}_i}\T{w}\cdot\left[\mu\T{j}+\rot\left(\T{b_r}\right]\right)dv+\int_{\partial\TTTT{D}_i}\T{w}\cdot \mu\T{j_s} ds;
\end{align}
\item $\forall \psi \in H^1(\TTTT{D})$:
\begin{align}
\label{eq:3weak3_b}
\int_{\TTTT{D}_i}\grad \left(\psi\right) \cdot \T{a}  dv = 0.
\end{align}
\end{itemize}

Questa formulazione, per quanto efficacie ed utilizzabile, comporta la necessità dell'introduzione del campo scalare $\chi$ aggiuntivo. Oltre all'ovvio aumento delle dimensioni del sistema, questa formulazione può portare a cattivo condizionamento dello stesso e ad una lenta convergenza nell'utilizzo di solutori lineari. Viene qui  pertanto introdotta una formulazione mista \textit{regolarizzata} tramite il parametro scalare $k$

$\forall \T{w} \in H(\rot,\TTTT{D})$: 
\begin{align}
\label{eq:3weak5}
\int_{\TTTT{D}_i}\rot\left(\T{w}\right)\cdot \rot\left(\T{a}\right) dv +k\int_{\TTTT{D}_i} \T{w} \cdot \T{a} dv
= \nonumber\\
=\int_{\TTTT{D}_i}\T{w}\cdot\left[\mu\T{j}+\rot\left(\T{b_r}\right)\right]dv+\int_{\partial\TTTT{D}_i}\T{w}\cdot \mu\T{j_s} ds.
\end{align}

Può essere dimostrato che tale formulazione converge ad una soluzione unica $\T{a_k}$ e che tale soluzione coincide con quella del problema non regolarizzato per $k\rightarrow 0$. 

\'E utile esprimere il parametro $k$ in relazione a $\mu_0$. In relazione al condizionamento del problema i seguenti valori sono stati sperimentati con successo: 

\begin{equation}
\label{eq:3 regpar}
k = c \mu_0 , \qquad 1<= c <=1000
\end{equation}
\section{Forma debole per il problema elettrico}

In analogia con quanto appena trattato per il potenziale magnetico, si imposta qui una struttura per la costruzione del problema in forma debole per il problema elettrico.

 In Sezione \ref{subsec:1potenziale} si è visto come il campo elettrico possa essere scomposto nell'ambito del regime \textit{MQS} in una parte solenoidale $\T{e_f}$ e in una irrotazionale $\T{e_c}$. Confrontando le equazioni \eqref{eq:2e_coulomb}, \eqref{eq:2e_faraday} e \eqref{eq:2max12s} si può osservare come nel regime stazionario (elettrostatica) la parte solenoidale sia identicamente nulla. 

La determinazione di $\T{e_f}$ deriva dall'analisi magnetica trattata nella sezione precedente. 

Per ottenere $\T{e_c}$ occorre cercare una forma debole della prima delle equazioni \eqref{eq:3potqss}
unitamente alle condizioni di interfaccia:

\begin{equation}
\left\{
\begin{aligned}
&\Bigl\langle \grad\left(\phi\right)\times \T{n}\Bigr\rangle = 0 \\
&\Bigl\langle \grad\left(\phi\right)\cdot \T{n}\Bigr\rangle = \rho_s.
\end{aligned}
\right.
\label{eq:3interface2_a}
\end{equation}

Sia ora $\phi \in H^1\left(\TTTT{D}\right)$; integrando l'Eq. \eqref{eq:3potqss} e la seconda delle condizioni al contorno \eqref{eq:3interface2_a} e moltiplicandole per una \textit{test function} $w \in H^1\left(\TTTT{D}\right)$ si ha :

\begin{align}
\label{eq:3weak_e1}
&-\int_{\TTTT{D}}w \, \lap\left(\phi\right) dv+\int_{\partial\TTTT{D}}w \, \grad\left(\phi\right)\cdot \T{n}ds= \nonumber\\
&=\varepsilon^{-1}\int_{\TTTT{D}}w\,\rho dv+\varepsilon^{-1}\int_{\partial\TTTT{D}}w\, \rho_s ds
\end{align}

\begin{eqnarray}
\label{eq:3weak_e2}
-\int_{\TTTT{D}}\dive\left[w \, \grad\left(\phi\right)\right]dv+ \int_{\TTTT{D}}\grad\left(w\right)\cdot \grad\left(\phi\right)dv+\\ +\int_{\partial\TTTT{D}}w \, \grad\left(\phi\right)\cdot \T{n}ds= \nonumber\\
=\varepsilon^{-1}\int_{\TTTT{D}}w\,\rho dv+\varepsilon^{-1}\int_{\partial\TTTT{D}}w\, \rho_s ds.
\end{eqnarray}

Utilizzando il teorema di Gauss per il primo termine, si ha:
\begin{eqnarray}
\label{eq:3weak_e3}
-\int_{\partial\TTTT{D}}w \, \grad\left(\phi\right)ds +\int_{\TTTT{D}}\grad\left(w\right)\cdot \grad\left(\phi\right)dv+\nonumber\\ +\int_{\partial\TTTT{D}}w \, \grad\left(\phi\right)\cdot \T{n}ds= \nonumber\\
=\varepsilon^{-1}\int_{\TTTT{D}}w\,\rho dv+\varepsilon^{-1}\int_{\partial\TTTT{D}}w\, \rho_s ds,
\end{eqnarray}
da cui
\begin{align}
\label{eq:3weak_e4}
&\int_{\TTTT{D}}\grad\left(w\right)\cdot \grad\left(\phi\right)dv = \nonumber\\
&=\varepsilon^{-1}\int_{\TTTT{D}}w\,\rho dv+\varepsilon^{-1}\int_{\partial\TTTT{D}}w\, \rho_s ds.
\end{align}

Come per il problema magnetico, per $w$ e $\phi \in H^1\left(\TTTT{D}\right)$ è dimostrabile l'esistenza e unicità della soluzione. \'E inoltre implicitamente soddisfatta la prima delle \eqref{eq:3interface2_a} \citep{zaglmayr_high_2006}.

\section{Conclusioni}

Nelle precedenti sezioni sono state presentate le forme deboli utilizzate nel corso di questo lavoro. 

L'esistenza e unicità delle soluzioni dei precedenti problemi sono garantite previa corretta impostazione del problema: il codice a elementi finiti deve pertanto utilizzare elementi conformi agli spazi funzionali presenti. I classici elementi finiti "lagrangiani" con  incognite nodali sono conformi agli spazi $H^1(\TTTT{D})$ e $H^(div, \TTTT{D})$, ma per la conformità ad $H(rot,\TTTT{D})$ si è reso necessario l'utilizzo di una diversa classe di elementi: gli elementi di spigolo (\textit{edge elements}). La principale proprietà di questi elementi è la possibilità di ottenere la formulazione di un problema la cui soluzione è richiesta conforme ad $H(rot,\TTTT{D})$ in termini di gradi di libertà esplicitamente indipendenti \citep{bossavit_computational_1998}. Fra gli elementi appartenenti a questa classe si è scelto di utilizzare gli elementi di Nédélec \citep{nedelec_acoustic_2001}. L'utilizzo di tali elementi, rispetto ad elementi nodali comporta grossi vantaggi per i problemi in forma debole presentati:
\begin{itemize}
\item Maggiore precisione a parità di griglia di calcolo
\item Miglior condizionamento della matrice assemblata
\end{itemize}

Ne risulta una robusta convergenza della soluzione con metodi iterativi.


In Tabella \ref{tab:tabella_campi} sono riassunti i risultati ottenuti in uno schema che mostra, per ciascun campo incognito, lo spazio funzionale adeguato e la classe di elementi utilizzata in questa analisi.

\begin{table}[h]
\centering
\caption{Elenco degli spazi funzionali appropriati e degli elementi utilizzati per i campi in gioco}
\footnotesize
\label{tab:tabella_campi}
\begin{tabular}{lcccc}
\toprule
Campo & Simbolo &Tipologia & Spazio funzionale & Elemento\\
\midrule
Campo elettrico & $\T{e}$ & vettoriale & $H(rot,\TTTT{D})$ & Nèdèlec\\
Campo magnetico & $\T{h}$ & vettoriale & $H(rot,\TTTT{D})$ & Nèdèlec\\
Induzione elettrica & $\T{d}$ & vettoriale & $H(div,\TTTT{D})$ & Lagrange\\
Induzione magnetica & $\T{b}$ & vettoriale & $H(div,\TTTT{D})$ & Lagrange\\
Potenziale elettrico & $\phi$ & vettoriale & $H^1(\TTTT{D})$ & Lagrange\\
Potenziale magnetico & $\T{a}$ & vettoriale & $H(rot,\TTTT{D})$ & Nèdèlec\\
\bottomrule
\end{tabular}
\end{table}

\section{Scaling}
I calcoli sono stati eseguiti in alcuni casi in parallelo sfruttando diverse \textit{CPU} per l'assemblaggio e la risoluzione del sistema lineare. Si riportano in questa sezione i risultati dei test di scalabilità effettuati per la valutazione delle prestazioni del solutore in funzione del numero di \textit{core} utilizzati.

\subsection{Strong scaling}
Per \textit{strong scaling} si intende l'analisi del tempo di calcolo necessario al programma per risolvere un determinato sistema al variare del numero di core. In Figura \ref{fig:strongscaling} è riportato il grafico delle prestazioni ottenute su un problema di esempio di calcolo del campo magnetico attorno a due magneti permanenti. Il grafico mostra soltanto i tempi di calcolo in parallelo, che nel programma sono preceduti dalle operazioni di \textit{preprocessing} eseguite in seriale con un processore. In Tabella \ref{tab:strongscaling} sono elencati i dati del modello e i tempi ottenuti.

\begin{table}[p]
\centering
\begin{minipage}{0.45\textwidth}
\centering
\tiny
\begin{tikzpicture}
    \begin{axis}[axis y line = left,
                 axis x line = left,
    		     width = \textwidth,
    		     height = .3\textheight,
                 xlabel = {cores}, 
                 ylabel = {tempo (s)}]
                 \addplot[mark = *]
    file {./Immagini/BVP/strongscaling.dat};
    \end{axis}
\end{tikzpicture}
\captionof{figure}{\textit{Strong scaling}}
\label{fig:strongscaling}
\end{minipage}
\begin{minipage}{0.45\textwidth}
\centering
\tiny
\caption{Dati del modello e tempi ottenuti per l'analisi di \textit{strong scaling}}
\label{tab:strongscaling}
\begin{tabular}{lc}
\toprule
Dati del modello&\\
\midrule
N. di vertici & 97609\\
N. di celle & 576435\\
\midrule
Tempi di calcolo seriale\\
\midrule
Meshing (gmsh) & 60.9 s\\
Conversione mesh & 28.9 s\\
Preprocessing & 8.7 s\\
\midrule
Sistemi risolti&\\
Sistema 1 & $675552\times675552$\\
Sistema 2 & $2319483\times2319483$\\
Sistema 3 & $878481\times878481$\\
\midrule
Tempi di calcolo\\
\midrule
1 \textit{core} & 231.3 s\\
2 \textit{core} & 215.7 s\\
4 \textit{core} & 153.3 s\\
6 \textit{core} & 140.8 s\\
\bottomrule
\end{tabular}
\end{minipage}
\end{table}


\subsection{Weak scaling}
Per \textit{weak scaling} si intende l'analisi delle prestazioni misurando il tempo di calcolo al variare del numero di core assegnando a ciascun processore la stessa mole di calcolo (quindi raddoppiando le dimensioni del sistema quando raddoppia il numero di core). In Figura \ref{fig:weakscaling} è riportato il grafico delle prestazioni ottenute. In Tabella \ref{tab:weakscaling} sono elencati i dati del modello e i tempi ottenuti.
\begin{table}[p]
\centering
\begin{minipage}{0.45\textwidth}
\centering
\tiny
\begin{tikzpicture}
    \begin{axis}[axis y line = left,
                 axis x line = left,
    		     width = \textwidth,
    		     height = .3\textheight,
                 xlabel = {cores}, 
                 ylabel = {tempo (s)}]
                 \addplot[mark = *]
    file {./Immagini/BVP/weakscaling.dat};
    \end{axis}
\end{tikzpicture}
\captionof{figure}{\textit{Weak scaling}}
\label{fig:weakscaling}

\end{minipage}
\begin{minipage}{0.45\textwidth}
\centering
\tiny
\begin{tabular}{lc}
\toprule
1 Processore&\\
\midrule
N. di vertici & 5243\\
N. di celle & 30713\\
Sistemi risolti&\\
Sistema 1 & $36066\times36066$\\
Sistema 2 & $123927\times123927$\\
Sistema 3 & $47187\times47187$\\
Tempo di calcolo & 10.0 s\\
\midrule
2 Processori&\\
\midrule
N. di vertici & 9686\\
N. di celle & 56970\\
Sistema 1 & $66827\times66827$\\
Sistema 2 & $229539\times229539$\\
Sistema 3 & $87174\times87174$\\
Tempo di calcolo & 13.6862208843 s\\
\midrule
4 Processori&\\
\midrule
N. di vertici & 18757\\
N. di celle & 110294\\
Sistemi risolti&\\
Sistema 1 & $129434\times129434$\\
Sistema 2 & $444573\times444573$\\
Sistema 3 & $444573\times444573$\\
Tempo di calcolo & 25.1 s\\

\bottomrule
\end{tabular}
\caption{Dati del modello e tempi ottenuti per l'analisi di \textit{weak scaling}}
\footnotesize
\label{tab:weakscaling}
\end{minipage}
\end{table}

I risultati ottenuti non sono soddisfacenti; si ritiene possa essere opportuno estendere il lavoro sviluppando un'interfaccia per il solutore \textit{AMS} dell libreria \textit{hypre}, che rappresenta il primo solutore scalabile per problemi elettromagnetici \citep{AMS_hypre}.